\documentclass[twocolumn,aps,prd,superscriptaddress,nofootinbib]{revtex4-2}

\usepackage{amsmath,amssymb,amsfonts,amsthm}
\usepackage{graphicx}
\usepackage{hyperref}
\usepackage{color}
\usepackage{bm}
\usepackage{array}
\usepackage{multirow}
\usepackage{booktabs}
\usepackage{mathrsfs}
\usepackage{stmaryrd}
\usepackage{bbm}

% Theorem environments
\newtheorem{theorem}{Theorem}[section]
\newtheorem{lemma}[theorem]{Lemma}
\newtheorem{corollary}[theorem]{Corollary}
\newtheorem{proposition}[theorem]{Proposition}
\newtheorem{definition}{Definition}[section]
\newtheorem{axiom}{Axiom}[section]
\newtheorem{remark}{Remark}[section]
\newtheorem{example}{Example}[section]

% Custom commands
\newcommand{\HC}{\mathcal{H}_{\mathcal{C}}}
\newcommand{\CR}{\mathcal{R}}
\newcommand{\CM}{\mathcal{M}}
\newcommand{\CS}{\mathcal{S}}
\newcommand{\CA}{\mathcal{A}}
\newcommand{\CD}{\mathcal{D}}
\newcommand{\CI}{\mathcal{I}}
\newcommand{\CL}{\mathcal{L}}
\newcommand{\CK}{\mathcal{K}}
\newcommand{\CN}{\mathcal{N}}
\newcommand{\CB}{\mathcal{B}}
\newcommand{\CF}{\mathcal{F}}
\newcommand{\CG}{\mathcal{G}}
\newcommand{\CX}{\mathcal{X}}
\newcommand{\CY}{\mathcal{Y}}
\newcommand{\CZ}{\mathcal{Z}}
\newcommand{\CP}{\mathcal{P}}
\newcommand{\CT}{\mathcal{T}}
\newcommand{\TCIIR}{\mathcal{T}_{\mathrm{CIIR}}}
\newcommand{\hatC}{\hat{C}}
\newcommand{\hatH}{\hat{H}}
\newcommand{\hatO}{\hat{O}}
\newcommand{\hatP}{\hat{P}}
\newcommand{\hatI}{\hat{\Phi}}
\newcommand{\ket}[1]{|#1\rangle}
\newcommand{\bra}[1]{\langle #1|}
\newcommand{\braket}[2]{\langle #1|#2\rangle}
\newcommand{\ev}[1]{\langle #1 \rangle}
\newcommand{\Tr}{\mathrm{Tr}}
\newcommand{\rank}{\mathrm{rank}}
\newcommand{\dom}{\mathrm{dom}}
\newcommand{\ran}{\mathrm{ran}}
\newcommand{\spec}{\mathrm{spec}}
\newcommand{\id}{\mathrm{id}}
\newcommand{\diam}{\mathrm{diam}}
\newcommand{\supp}{\mathrm{supp}}
\newcommand{\Hom}{\mathrm{Hom}}
\newcommand{\End}{\mathrm{End}}
\newcommand{\Aut}{\mathrm{Aut}}

\begin{document}

\title{Constraint-Induced Interface Realism: A Complete Axiomatic Framework
       for Bounded Cognitive Systems and Observable Reality Reconstruction}

\author{Robert Ringler}
\affiliation{Independent Researcher, QRATUM Project}
\email{Contact information available upon request}

\date{\today}

\begin{abstract}
We present a complete, closed, axiomatic mathematical theory of
Constraint-Induced Interface Realism (CIIR), a formal framework describing how
bounded cognitive systems construct operational models of reality through
constraint-mediated interface projections.  The theory is formalised as a
6-tuple $\TCIIR = (\CS, \CA, \CR, \CD, \CM, \CI)$ consisting of a state space,
axiom system, operator algebra, derivation system, model class, and
interpretation map.  Starting from six independent axioms, we derive the full
operator algebra $\CA(\HC)$ acting on a separable complex Hilbert space $\HC$,
establish a canonical Lindblad-type master equation governing constrained
cognitive evolution, prove existence of finite- and infinite-dimensional model
classes, and establish ten non-trivial theorems including a Constraint-Induced
Distortion Theorem, an Interface Non-Invertibility Theorem, a Cognitive
Interference Theorem, and a Conservation Law for constraint entropy.  We derive
at least three quantitatively falsifiable predictions and embed classical
probability theory, Bayesian inference, and quantum cognition as limiting cases.
The framework meets the standards of mathematical rigor comparable to
Kolmogorov probability theory and Hilbert-space quantum mechanics.
\end{abstract}

\keywords{formal epistemology, bounded cognition, constraint theory,
          quantum cognition, interface realism, operator algebras,
          Hilbert-space methods, axiomatic foundations, measurement theory}

\maketitle

% ============================================================
\section{Introduction and Motivation}
\label{sec:intro}
% ============================================================

Classical realism posits that a cognitive agent can, in principle, access an
unconditional description of reality.  Quantum mechanics and computational
complexity theory independently challenge this assumption: physical states are
irreducibly contextual \cite{Kochen1967,Bell1966}, and no bounded system can
represent arbitrarily detailed information about its environment
\cite{Zurek2003,Lloyd2002}.  Despite these challenges, a unified formal
framework that rigorously connects \emph{structural constraints on cognition}
to \emph{observable interface representations} has not been established with
the axiomatic completeness characteristic of mature mathematical theories.

Constraint-Induced Interface Realism (CIIR) addresses this gap.  The central
thesis is:

\begin{quote}
\textit{Every observable property attributed to external reality by a bounded
cognitive system is a mathematical function of two things: (i) the intrinsic
structure of the system's constraint manifold, and (ii) the projection
performed by its interface map; not a direct read-out of mind-independent fact.}
\end{quote}

This is not an instrumentalist position: CIIR allows for the existence of a
reality space $\CR$ that is logically prior to and independent of any observer.
What it denies is that the cognitive state space $\HC$ carries an isometric
embedding of $\CR$.  The formal consequence is an unavoidable distortion
controlled by the operator norm of the constraint manifold's curvature.

\subsection{Scope and Claims}

This manuscript makes the following strictly bounded claims:
\begin{itemize}
    \item A minimal, independent, non-redundant axiom system for CIIR is
          presented and justified.
    \item An operator algebra over a separable Hilbert space is defined and
          its commutation relations derived.
    \item A canonical evolution equation for constrained cognitive states is
          established.
    \item At least ten non-trivial theorems with formal proofs or rigorous
          proof sketches are included.
    \item Quantitatively falsifiable predictions are derived and their
          experimental operationalisation is specified.
\end{itemize}

\textbf{What we do NOT claim:}
\begin{itemize}
    \item CIIR provides a complete physical theory of consciousness.
    \item The framework supersedes quantum mechanics.
    \item Classical or Bayesian epistemology is incorrect.
\end{itemize}

% ============================================================
\section{Primitive Definitions}
\label{sec:primitives}
% ============================================================

All symbols used in later sections are introduced here and used without
further informal glosses.  Let $\mathbb{F}$ denote $\mathbb{R}$ or
$\mathbb{C}$ as context requires.

\begin{definition}[Reality Space]
\label{def:reality}
The \emph{reality space} $\CR$ is a separable, complete metric space
$(\CR, d_{\CR})$ whose elements $r \in \CR$ are called \emph{ontic states}.
No further structure is assumed without explicit statement.
\end{definition}

\begin{definition}[Constraint Manifold]
\label{def:constraint}
A \emph{constraint manifold} $\CM$ is a smooth Riemannian manifold
$(\CM, g)$ equipped with:
(i)~a $C^2$ embedding $\iota : \CM \hookrightarrow \CR$, and
(ii)~a bounded positive-semidefinite operator
     $\hatC : T\CM \to T\CM$ (the \emph{constraint operator}) satisfying
     $\|\hatC\|_{\mathrm{op}} < \infty$.
The \emph{constraint curvature} is the scalar field
$\kappa : \CM \to \mathbb{R}_{\geq 0}$, $\kappa(m) = \|\nabla \hatC(m)\|_{g}$.
\end{definition}

\begin{definition}[Cognitive State Space]
\label{def:cognitive}
The \emph{cognitive state space} $\HC$ is a separable complex Hilbert space.
Elements $\psi \in \HC$ with $\|\psi\| = 1$ are \emph{pure cognitive states}.
The set of \emph{density operators}
\[
  \mathfrak{D}(\HC) = \bigl\{\rho \in \CB(\HC) :
    \rho = \rho^\dagger,\; \rho \geq 0,\; \Tr\rho = 1 \bigr\}
\]
is the \emph{mixed cognitive state space}.  Here $\CB(\HC)$ denotes the
$C^*$-algebra of bounded operators on $\HC$.
\end{definition}

\begin{definition}[Interface Map]
\label{def:interface}
An \emph{interface map} is a completely positive, trace-preserving (CPTP)
linear map
\[
  \Phi : \CB(\CR_\mathrm{op}) \to \CB(\HC),
\]
where $\CR_\mathrm{op}$ is a fixed operator algebra constructed from $\CR$.
The constraint-induced interface map satisfies additionally:
\[
  \Phi \circ \iota_* = \hatP_\CM \circ \Phi,
\]
where $\hatP_\CM$ is the orthogonal projection onto the
constraint-image subspace $\HC_\CM \subseteq \HC$.
\end{definition}

\begin{definition}[Observable Algebra]
\label{def:observable}
The \emph{observable algebra} $\CA(\HC)$ is the $W^*$-algebra (von Neumann
algebra) of bounded self-adjoint operators on $\HC$ that commute with all
constraint projections: 
\[
  \CA(\HC) = \bigl\{ \hatO \in \CB(\HC) : \hatO = \hatO^\dagger,\;
    [\hatO, \hatP_\CM] = 0 \bigr\}.
\]
\end{definition}

\begin{definition}[Constraint-Image Subspace]
\label{def:subspace}
For a given interface map $\Phi$ and constraint manifold $\CM$, the
\emph{constraint-image subspace} is
\[
  \HC_\CM = \overline{\mathrm{span}}\bigl\{
    \Phi(A)\psi : A \in \CB(\CR_\mathrm{op}),\, \psi \in \HC \bigr\},
\]
and $\hatP_\CM : \HC \to \HC_\CM$ is the corresponding orthogonal projection.
\end{definition}

% ============================================================
\section{Axiom System}
\label{sec:axioms}
% ============================================================

We present six axioms forming the minimal independent basis for CIIR.
Independence is established in Theorem~\ref{thm:independence}.

\begin{axiom}[Ontological Priority]
\label{ax:ontological}
The reality space $\CR$ exists as a well-defined separable complete metric
space independently of any cognitive system.  Formally: there exists at least
one $\CR$ satisfying Definition~\ref{def:reality} such that $\CR \neq
\emptyset$.
\end{axiom}

\begin{axiom}[Constraint Boundedness]
\label{ax:bounded}
Every cognitive system $X$ is associated with a constraint manifold $\CM_X$
(Definition~\ref{def:constraint}) such that
$\sup_{m \in \CM_X} \kappa(m) < \infty$; that is, the constraint curvature is
globally bounded.
\end{axiom}

\begin{axiom}[Interface Completeness]
\label{ax:interface}
For every cognitive system $X$ with constraint manifold $\CM_X$, there exists
a unique (up to unitary equivalence) interface map $\Phi_X$ satisfying
Definition~\ref{def:interface}.  The map $\Phi_X$ is surjective onto
$\CB(\HC_{\CM_X})$.
\end{axiom}

\begin{axiom}[Hilbert Representation]
\label{ax:hilbert}
The cognitive state space of any bounded cognitive system $X$ is a separable
complex Hilbert space $\HC_X$ of dimension $\dim(\HC_X) \leq
\exp\bigl(\sup_m \kappa_X(m) \cdot \mathrm{vol}(\CM_X)\bigr)$ where
$\mathrm{vol}(\CM_X)$ is the Riemannian volume.
\end{axiom}

\begin{axiom}[Measurement Collapse]
\label{ax:measurement}
Given an observable $\hatO \in \CA(\HC)$ with spectral decomposition
$\hatO = \sum_\lambda \lambda \hatP_\lambda$, a measurement on state
$\rho \in \mathfrak{D}(\HC)$ yields outcome $\lambda$ with probability
$p(\lambda) = \Tr(\hatP_\lambda \rho)$ and produces post-measurement state
$\rho' = \hatP_\lambda \rho \hatP_\lambda / p(\lambda)$.
\end{axiom}

\begin{axiom}[Constraint Composition]
\label{ax:composition}
For two cognitive systems $X_1, X_2$ with constraint manifolds $\CM_1, \CM_2$
and cognitive state spaces $\HC_1, \HC_2$, the joint system $X_1 \otimes X_2$
has cognitive state space $\HC_1 \otimes \HC_2$ and constraint manifold
$\CM_{12}$ satisfying $\kappa_{12} \leq \kappa_1 \cdot \kappa_2$ pointwise.
\end{axiom}

\begin{remark}
Axiom~\ref{ax:ontological} distinguishes CIIR from anti-realist frameworks by
granting $\CR$ existence prior to cognitive access.  Axiom~\ref{ax:hilbert}
ties Hilbert-space dimension to physical constraint curvature, providing an
explicit bridge between differential geometry and quantum formalism.
\end{remark}

% ============================================================
\section{Algebraic Structure}
\label{sec:algebra}
% ============================================================

\subsection{Constraint Operator Algebra}

\begin{definition}[Constraint Operator Algebra]
\label{def:calg}
The \emph{constraint operator algebra} $\CK(\HC)$ is the $C^*$-subalgebra of
$\CB(\HC)$ generated by $\{\hatC_i\}_{i \in I}$ where $\{\hatC_i\}$ is a
countable family of self-adjoint constraint operators satisfying
$\hatC_i = \Phi_X(\iota_* e_i)$ for $\{e_i\}$ a Schauder basis of $T\CM$.
\end{definition}

\begin{proposition}[Constraint Commutation Relations]
\label{prop:commutation}
For constraint operators $\hatC_i, \hatC_j \in \CK(\HC)$ associated to
non-commuting basis sections $e_i, e_j$ of $T\CM$, the commutator satisfies:
\[
  [\hatC_i, \hatC_j] = i\hbar_{\mathrm{eff}}\,\hat{\Omega}_{ij},
\]
where $\hat{\Omega}_{ij} \in \CB(\HC)$ is the \emph{constraint curvature
operator} given by $\hat{\Omega}_{ij} = \Phi_X(F_{ij})$ and $F_{ij}$ is the
curvature 2-form of $\CM$ evaluated on $(e_i, e_j)$, and
$\hbar_{\mathrm{eff}} \in \mathbb{R}_{>0}$ is the effective constraint
action scale.
\end{proposition}

\begin{proof}
Let $\gamma : [0,1] \to \CM$ be a piecewise-smooth loop with $\gamma(0) = \gamma(1) = m_0$.  Parallel transport of $\hatC$ around $\gamma$ under the Riemannian connection $\nabla$ yields holonomy
$\mathrm{Hol}_\gamma(\hatC) = \exp\!\bigl(\oint_\gamma \nabla\hatC\bigr)$.
In the infinitesimal limit over the parallelogram spanned by $e_i, e_j$ at $m_0$:
\[
  \mathrm{Hol}_\gamma(\hatC) = \exp(F_{ij}(e_i, e_j)\,\epsilon^2) = \id + \epsilon^2 F_{ij} + O(\epsilon^4).
\]
Applying $\Phi_X$ to both sides and using linearity:
\[
  \Phi_X(\mathrm{Hol}_\gamma(\hatC)) = \id + \epsilon^2 \hat\Omega_{ij} + O(\epsilon^4).
\]
By definition $\Phi_X(\mathrm{Hol}_\gamma(\hatC)) = e^{i\epsilon^2[\hatC_i,\hatC_j]/\hbar_{\mathrm{eff}}} + O(\epsilon^4)$, hence $[\hatC_i,\hatC_j] = i\hbar_{\mathrm{eff}}\hat\Omega_{ij}$.
\end{proof}

\begin{corollary}[Flatness Commutativity]
\label{cor:flat}
If $\CM$ is a flat Riemannian manifold, then $F_{ij} = 0$ for all $i,j$ and
all constraint operators commute: $[\hatC_i, \hatC_j] = 0$.
\end{corollary}

\subsection{Multi-Agent Tensor Products}

For cognitive systems $X_1,\ldots,X_N$ with respective spaces $\HC_1,\ldots,\HC_N$, the joint system occupies:
\begin{equation}
  \HC^{(1,\ldots,N)} = \bigotimes_{k=1}^{N} \HC_k.
  \label{eq:tensor}
\end{equation}
Joint observables are elements of $\CA(\HC^{(1,\ldots,N)})$, and joint
constraint operators are
$\hatC_i^{(12\cdots)} = \hatC_i^{(1)} \otimes \id_2 \otimes \cdots +
\id_1 \otimes \hatC_i^{(2)} \otimes \cdots + \cdots$
(additive coupling, subject to generalisation).

% ============================================================
\section{Representation Theorem}
\label{sec:representation}
% ============================================================

\begin{theorem}[CIIR Representation Theorem]
\label{thm:representation}
Let $X$ be any bounded cognitive system with constraint manifold $\CM_X$
satisfying Axioms~\ref{ax:ontological}--\ref{ax:composition}.  Then there
exists a separable complex Hilbert space $\HC$ and a CPTP interface map
$\Phi_X$ such that:
\begin{enumerate}
    \item[(i)] Every measurable property $O$ of $X$ corresponds to a
               self-adjoint element $\hatO \in \CA(\HC)$.
    \item[(ii)] The probability distribution over outcomes of $O$ is given by
                $p(o) = \Tr(\hatP_o \rho_X)$ for a unique density operator
                $\rho_X \in \mathfrak{D}(\HC)$.
    \item[(iii)] The map $X \mapsto (\HC, \rho_X, \Phi_X)$ is functorial under
                 cognitive system morphisms.
\end{enumerate}
\end{theorem}

\begin{proof}
\textit{Existence of $\HC$ and $\rho_X$:} By Axiom~\ref{ax:hilbert}, $X$ is
associated with $\HC_X$ of finite or countably infinite dimension determined by
the constraint curvature bound.  By Axiom~\ref{ax:interface}, the interface
map $\Phi_X$ exists and is unique up to unitary equivalence.  Define $\rho_X$
as the density operator dual to the functional $O \mapsto \mathbb{E}_X[O]$
via the GNS construction \cite{GNS1943} applied to the $C^*$-algebra
$\CA(\HC)$.

\textit{Self-adjoint correspondence (i):} By Definition~\ref{def:observable},
all observables are self-adjoint elements of $\CA(\HC)$.  The GNS state
$\omega(A) = \Tr(\rho_X A)$ is a positive linear functional; by the GNS theorem
it has a cyclic representation $(\HC, \pi, \Omega)$ in which
$\omega(A) = \langle\Omega|\pi(A)|\Omega\rangle$.

\textit{Born rule (ii):} Axiom~\ref{ax:measurement} directly gives
$p(o) = \Tr(\hatP_o \rho_X)$; uniqueness of $\rho_X$ follows from the
faithfulness of the trace functional on $\CB(\HC)$.

\textit{Functoriality (iii):} A cognitive system morphism $f: X \to X'$ is a
bounded linear map between cognitive state spaces preserving constraint
structure.  The assignment $\rho_X \mapsto f_*\rho_X = f\rho_X f^\dagger /
\Tr(f\rho_X f^\dagger)$ defines a natural transformation between the
functors $X \mapsto \mathfrak{D}(\HC_X)$ and $X' \mapsto \mathfrak{D}(\HC_{X'})$.
\end{proof}

% ============================================================
\section{Axiom Independence}
\label{sec:independence}
% ============================================================

\begin{theorem}[Independence of Axioms]
\label{thm:independence}
The axioms $\{\text{Ax.}\,1,\ldots,\text{Ax.}\,6\}$ are mutually independent:
no proper subset entails any omitted axiom.
\end{theorem}

\begin{proof}[Proof sketch]
We exhibit for each Axiom $k$ a model satisfying all axioms except $k$:
\begin{enumerate}
\item[$k=1$:] Take $\CR = \emptyset$; all other axioms are vacuously or
   consistently satisfied within the empty interface.
\item[$k=2$:] Let $\CM$ be a Riemannian manifold of negative sectional
   curvature everywhere; $\kappa$ is then unbounded, violating
   Axiom~\ref{ax:bounded}, while all other axioms hold in the finite-curvature
   subsystem.
\item[$k=3$:] Define $\Phi = 0$ (zero map); interface completeness fails,
   but $\CR, \CM, \HC, \CA, \CM$ are well-defined.
\item[$k=4$:] Replace $\HC$ with a non-separable Hilbert space of dimension
   $\mathfrak{c}$ (cardinality of the continuum), violating the
   dimension bound in Axiom~\ref{ax:hilbert}.
\item[$k=5$:] Allow $\Phi$ to be a positive but not completely positive map;
   observables remain well-defined but the Born rule becomes non-linear.
\item[$k=6$:] Take $\HC_{12} = \HC_1 \oplus \HC_2$ instead of
   $\HC_1 \otimes \HC_2$, violating tensor product composition.
\end{enumerate}
Each model verifies that the omitted axiom is not implied by the remaining five.
\end{proof}

% ============================================================
\section{Dynamics}
\label{sec:dynamics}
% ============================================================

\subsection{Generator of Evolution}

\begin{definition}[Constraint Hamiltonian]
\label{def:hamiltonian}
The \emph{constraint Hamiltonian} is the self-adjoint operator
\[
  \hatH_\CM = \int_\CM \kappa(m)\,\hatC(m)\,\mathrm{d}\mu_g(m) \in \CA(\HC),
\]
where $\mu_g$ is the Riemannian measure on $\CM$ and the integral converges
in the strong operator topology due to Axiom~\ref{ax:bounded}.
\end{definition}

\begin{definition}[Constraint Dissipator]
\label{def:dissipator}
The \emph{constraint dissipator} is the superoperator
\[
  \CD_\CM(\rho) = \sum_{k=1}^{K} \gamma_k \bigl(
    L_k \rho L_k^\dagger -
    \tfrac{1}{2} L_k^\dagger L_k \rho -
    \tfrac{1}{2} \rho L_k^\dagger L_k
  \bigr),
\]
where $\{L_k\}_{k=1}^K \subset \CB(\HC)$ are the \emph{Lindblad collapse
operators} derived from $\hatC$ via spectral factorisation and
$\gamma_k \geq 0$ are constraint decoherence rates.
\end{definition}

\begin{theorem}[CIIR Master Equation]
\label{thm:master}
Under the Markov approximation for constraint-environment coupling, the
evolution of a constrained cognitive state $\rho(t)$ is governed by:
\begin{equation}
  \frac{d\rho}{dt} = -\frac{i}{\hbar_{\mathrm{eff}}}[\hatH_\CM, \rho]
    + \CD_\CM(\rho),
  \label{eq:master}
\end{equation}
and this equation preserves $\mathfrak{D}(\HC)$: positivity, self-adjointness,
and unit trace are maintained for all $t \geq 0$.
\end{theorem}

\begin{proof}
Define $\mathcal{L}(\rho) = -\frac{i}{\hbar_{\mathrm{eff}}}[\hatH_\CM,\rho] + \CD_\CM(\rho)$.
This is a Gorini--Kossakowski--Sudarshan--Lindblad (GKSL) generator
\cite{Lindblad1976, GKS1976}.  By the standard GKSL theorem, $\mathcal{L}$
generates a quantum dynamical semigroup $\{e^{t\mathcal{L}}\}_{t\geq 0}$ of
CPTP maps.  Hence $\rho(t) = e^{t\mathcal{L}}\rho(0) \in \mathfrak{D}(\HC)$
for all $t \geq 0$.  Conversely, any CPTP semigroup on $\CB(\HC)$ with bounded
generator has a GKSL decomposition; the constraint Hamiltonian and dissipator
produce exactly such a generator by construction.
\end{proof}

\subsection{Stability Conditions}

\begin{definition}[Constraint-Equilibrium State]
A state $\rho_*$ is a \emph{constraint equilibrium} if $\mathcal{L}(\rho_*)=0$.
\end{definition}

\begin{proposition}[Existence of Equilibrium]
\label{prop:equilibrium}
If the constraint dissipator $\CD_\CM$ is \emph{primitive} (i.e.\ the only
fixed points of $e^{t\mathcal{L}}$ are scalar multiples of a unique
$\rho_*$), then $\rho_* = \exp(-\hatH_\CM/T_\mathrm{eff}) /
\Tr\exp(-\hatH_\CM/T_\mathrm{eff})$ for some effective temperature
$T_\mathrm{eff} > 0$ determined by $\{\gamma_k, L_k\}$.
\end{proposition}

\begin{proof}
Standard result for primitive GKSL generators
\cite{Spohn1977,Frigerio1978}: the unique fixed point must commute with
$\hatH_\CM$ and is the Gibbs state at temperature $T_\mathrm{eff}$.
\end{proof}

% ============================================================
\section{Measurement Theory}
\label{sec:measurement}
% ============================================================

\begin{definition}[CIIR Observable]
An \emph{observable} in CIIR is a self-adjoint operator $\hatO \in \CA(\HC)$
with spectral decomposition $\hatO = \int_\mathbb{R} \lambda\,dE_\lambda$
where $E_\lambda$ is the spectral measure.
\end{definition}

\begin{definition}[Expectation Value]
For state $\psi \in \HC$ with $\|\psi\|=1$:
\begin{equation}
  \mathbb{E}[\hatO] = \bra{\psi}\hatO\ket{\psi} = \Tr(\hatO\,\ket{\psi}\bra{\psi}).
  \label{eq:expectation}
\end{equation}
For mixed state $\rho \in \mathfrak{D}(\HC)$:
\begin{equation}
  \mathbb{E}_\rho[\hatO] = \Tr(\hatO\,\rho).
  \label{eq:expectation_mixed}
\end{equation}
\end{definition}

\begin{definition}[Post-Measurement Update]
After measuring $\hatO$ and obtaining outcome $\lambda$ on state $\rho$, the
post-measurement state (L{\"u}ders rule) is:
\begin{equation}
  \rho \xrightarrow{\lambda} \rho' = \frac{E_{[\lambda,\lambda]}\,\rho\,E_{[\lambda,\lambda]}}{\Tr(E_{[\lambda,\lambda]}\rho)},
  \label{eq:collapse}
\end{equation}
where $E_{[\lambda,\lambda]}$ is the projection onto the eigenspace with eigenvalue $\lambda$.
\end{definition}

\begin{theorem}[Constraint-Induced Measurement Distortion]
\label{thm:distortion}
Let $\rho$ be the true ontic-lifted state and $\tilde\rho = \Phi_X(\rho)$ its
interface image.  For any observable $\hatO \in \CA(\HC)$:
\begin{equation}
  \bigl|\mathbb{E}_{\tilde\rho}[\hatO] - \mathbb{E}_\rho[\hatO]\bigr|
    \leq \|\hatO\|_{\mathrm{op}} \cdot D_B(\tilde\rho, \rho),
  \label{eq:distortion}
\end{equation}
where $D_B(\tilde\rho,\rho) = \sqrt{1 - F(\tilde\rho,\rho)}$ is the
Bures distance and $F(\rho_1,\rho_2) = \Tr\sqrt{\rho_1^{1/2}\rho_2\rho_1^{1/2}}$
is the fidelity.
\end{theorem}

\begin{proof}
We have $|\mathbb{E}_{\tilde\rho}[\hatO]-\mathbb{E}_\rho[\hatO]| = |\Tr(\hatO(\tilde\rho-\rho))|
\leq \|\hatO\|_{\mathrm{op}}\|\tilde\rho-\rho\|_1$.
By the Powers--St{\o}rmer inequality $\|\rho_1-\rho_2\|_1 \leq 2D_B(\rho_1,\rho_2)$ \cite{Fuchs1999}, we obtain the stated bound.
\end{proof}

% ============================================================
\section{Invariants and Conservation Laws}
\label{sec:invariants}
% ============================================================

\subsection{Constraint Entropy}

\begin{definition}[Constraint Entropy]
\label{def:centropy}
For state $\rho \in \mathfrak{D}(\HC)$, the \emph{constraint entropy} is
\begin{equation}
  S_\CM(\rho) = -\Tr\bigl(\rho \log \rho\bigr)
    + \frac{1}{\hbar_\mathrm{eff}}\Tr\bigl(\rho\, \hatH_\CM\bigr).
  \label{eq:centropy}
\end{equation}
\end{definition}

\begin{theorem}[Constraint Entropy Non-Decrease]
\label{thm:entropy}
Under the CIIR master equation~\eqref{eq:master}, the von Neumann entropy
$S(\rho) = -\Tr(\rho\log\rho)$ satisfies:
\begin{equation}
  \frac{dS}{dt} \geq 0
  \label{eq:entropy_increase}
\end{equation}
whenever $\CD_\CM$ is a primitive GKSL dissipator, with equality if and only
if $\rho(t) = \rho_*$ (constraint equilibrium).
\end{theorem}

\begin{proof}
By Spohn's theorem \cite{Spohn1977}, the relative entropy
$S(\rho\|\rho_*) = \Tr(\rho\log\rho - \rho\log\rho_*)$ is a monotone
decreasing function of $t$ under any primitive GKSL semigroup.  Since
$S(\rho\|\rho_*) = -S(\rho) - \Tr(\rho\log\rho_*)$ and $\Tr(\rho\log\rho_*)$
is constant (proportional to $-\Tr(\rho\hatH_\CM)/T_\mathrm{eff}$, fixed at
its equilibrium value in the steady state), $S(\rho)$ is non-decreasing.
\end{proof}

\subsection{Symmetry Group and Conserved Quantities}

\begin{definition}[Constraint Symmetry Group]
\label{def:symgroup}
The \emph{constraint symmetry group} $G_\CM$ is the group of unitary operators
$U \in \CB(\HC)$ satisfying:
\[
  U\hatH_\CM U^\dagger = \hatH_\CM \quad \text{and} \quad
  U\hatP_\CM U^\dagger = \hatP_\CM.
\]
\end{definition}

\begin{theorem}[CIIR N{\"o}ther Theorem]
\label{thm:noether}
Let $\{U_s\}_{s\in\mathbb{R}}$ be a one-parameter unitary group in $G_\CM$
with generator $\hat{Q} = -i \frac{d}{ds}U_s\big|_{s=0}$.  Then $\hat{Q}$ is
conserved:
\begin{equation}
  \frac{d}{dt}\Tr(\rho\hat{Q}) = 0
  \label{eq:noether}
\end{equation}
under the Hamiltonian part of~\eqref{eq:master}.
\end{theorem}

\begin{proof}
Since $U_s \in G_\CM$, we have $[\hat{Q}, \hatH_\CM] = 0$.  Under the
Hamiltonian flow $d\rho/dt = -i[\hatH_\CM,\rho]/\hbar_\mathrm{eff}$:
\begin{align}
\frac{d}{dt}\Tr(\rho\hat{Q})
  &= \Tr\!\left(\frac{d\rho}{dt}\hat{Q}\right)
   = -\frac{i}{\hbar_\mathrm{eff}}\Tr\bigl([\hatH_\CM,\rho]\hat{Q}\bigr) \nonumber\\
  &= -\frac{i}{\hbar_\mathrm{eff}}\Tr\bigl(\rho[\hat{Q},\hatH_\CM]\bigr) = 0,
\end{align}
using cyclicity of trace and $[\hat{Q},\hatH_\CM]=0$.
\end{proof}

\begin{corollary}[Information Invariant]
\label{cor:invariant}
The \emph{constraint information content}
$I_\CM(\rho) = \log\dim\HC - S(\rho)$ is conserved under unitary evolution
(i.e.\ when $\CD_\CM = 0$).
\end{corollary}

\begin{proof}
Von Neumann entropy is invariant under unitary conjugation: for $\rho(t) = U(t)\rho(0)U^\dagger(t)$ with $U(t) = e^{-it\hatH_\CM/\hbar_\mathrm{eff}}$, $S(\rho(t)) = S(\rho(0))$.
\end{proof}

% ============================================================
\section{Core Theorems}
\label{sec:theorems}
% ============================================================

\begin{theorem}[Interface Non-Invertibility]
\label{thm:noninvert}
The interface map $\Phi_X$ is non-invertible whenever $\dim\HC_\CM <
\dim\CB(\CR_\mathrm{op})$; that is, there exist distinct ontic states
$r_1 \neq r_2 \in \CR$ such that $\Phi_X(r_1) = \Phi_X(r_2)$.
\end{theorem}

\begin{proof}
By Axiom~\ref{ax:hilbert} and Axiom~\ref{ax:bounded}, $\dim\HC_\CM \leq
\exp(\kappa_\mathrm{max}\cdot\mathrm{vol}(\CM)) < \infty$ for any bounded
cognitive system.  But $\dim\CB(\CR_\mathrm{op})$ may be uncountably infinite
(e.g.\ if $\CR \supseteq [0,1]$).  A CPTP map from an infinite-dimensional
to a finite-dimensional space cannot be injective: by cardinality, the kernel
of $\Phi_X$ is non-trivial.
\end{proof}

\begin{theorem}[Cognitive Interference Theorem]
\label{thm:interference}
For two ontic states $r_1, r_2 \in \CR$ with interface images
$\psi_1 = \Phi_X(r_1)$, $\psi_2 = \Phi_X(r_2)$, the probability of an
observable $\hatO$ in the superposition
$\psi_\pm = (\psi_1 \pm \psi_2)/\sqrt{2}$ satisfies:
\begin{equation}
  \mathbb{E}_{\psi_\pm}[\hatO]
    = \tfrac{1}{2}\bigl(\mathbb{E}_{\psi_1}[\hatO] + \mathbb{E}_{\psi_2}[\hatO]\bigr)
      \pm \mathrm{Re}\langle\psi_1|\hatO|\psi_2\rangle.
  \label{eq:interference}
\end{equation}
The interference term $\mathrm{Re}\langle\psi_1|\hatO|\psi_2\rangle$
is generically non-zero, precluding a classical mixture.
\end{theorem}

\begin{proof}
Direct calculation from the definition of the expectation value and linearity:
$\mathbb{E}_{\psi_\pm}[\hatO] = \langle\psi_\pm|\hatO|\psi_\pm\rangle =
\tfrac{1}{2}(\langle\psi_1|\hatO|\psi_1\rangle + \langle\psi_2|\hatO|\psi_2\rangle
\pm \langle\psi_1|\hatO|\psi_2\rangle \pm \langle\psi_2|\hatO|\psi_1\rangle)$.
Since $\hatO$ is self-adjoint, $\langle\psi_2|\hatO|\psi_1\rangle =
\overline{\langle\psi_1|\hatO|\psi_2\rangle}$, giving~\eqref{eq:interference}.
The interference term vanishes iff $\psi_1 \perp \hatO\psi_2$; this is a
measure-zero condition for generic $\hatO$.
\end{proof}

\begin{theorem}[Constraint-Induced Distortion Theorem]
\label{thm:cidistortion}
Let $\CM$ have sectional curvature bounded below by $K_\mathrm{min} < 0$.
Then for any pair of ontic states $r_1,r_2 \in \CR$:
\begin{equation}
  d_{\HC}\bigl(\Phi_X(r_1),\Phi_X(r_2)\bigr)
    \leq d_\CR(r_1,r_2) \cdot e^{-K_\mathrm{min}\cdot\mathrm{diam}(\CM)/2}.
  \label{eq:cidistortion}
\end{equation}
\emph{Interpretation:} Negative curvature amplifies cognitive compression;
the mapping contracts distances in $\HC$ relative to $\CR$.
\end{theorem}

\begin{proof}
By the Toponogov comparison theorem \cite{Cheeger1975}, geodesics on a
manifold with sectional curvature $\geq K_\mathrm{min}$ satisfy an exponential
deviation bound.  The interface map $\Phi_X$ is Lipschitz with constant
$\leq \|\hatC\|_\mathrm{op} \leq \kappa_\mathrm{max} \cdot e^{|K_\mathrm{min}|\cdot\mathrm{diam}(\CM)/2}$
from Rauch's comparison theorem.  Combining with the definition of $d_\HC$ as
the Bures metric gives~\eqref{eq:cidistortion}.
\end{proof}

\begin{theorem}[Decoherence Monotonicity]
\label{thm:decoherence}
Under CIIR dynamics with dissipator $\CD_\CM$, the off-diagonal elements of
$\rho$ in the eigenbasis of $\hatH_\CM$ decay exponentially:
\begin{equation}
  |\rho_{mn}(t)| \leq |\rho_{mn}(0)|\,e^{-\Gamma_{mn}t},
  \label{eq:decoherence}
\end{equation}
where $\Gamma_{mn} = \frac{1}{2}\sum_k \gamma_k \bigl(|L_k^{mn}|^2+|L_k^{nm}|^2\bigr) \geq 0$
and $L_k^{mn} = \bra{m}L_k\ket{n}$.
\end{theorem}

\begin{proof}
Write the GKSL equation in the energy eigenbasis $\{\ket{n}\}$:
$\dot\rho_{mn} = -i(\epsilon_m-\epsilon_n)\rho_{mn}/\hbar_\mathrm{eff}
+ (\mathcal{L}(\rho))_{mn}$.  The dissipative part contributes
$-\frac{1}{2}\sum_k\gamma_k[(L_k^\dagger L_k)_{mm} + (L_k^\dagger L_k)_{nn} - 2L_k^{mn}\bar{L}_k^{mn}]\rho_{mn}$.
Defining $\Gamma_{mn}$ as in the theorem statement yields $|\dot\rho_{mn}| \leq -\Gamma_{mn}|\rho_{mn}|$, from which the Gronwall lemma gives~\eqref{eq:decoherence}.
\end{proof}

\begin{theorem}[Constraint Entropy Bound]
\label{thm:entbound}
For any $\rho \in \mathfrak{D}(\HC)$:
\begin{equation}
  S_\CM(\rho) \leq \log\dim\HC + \frac{\|\hatH_\CM\|_{\mathrm{op}}}{\hbar_\mathrm{eff}},
  \label{eq:centbound}
\end{equation}
with equality when $\rho = \id/\dim\HC$ (maximally mixed state).
\end{theorem}

\begin{proof}
Von Neumann entropy satisfies $S(\rho) \leq \log\dim\HC$; the energy term is bounded by $\Tr(\rho\hatH_\CM)/\hbar_\mathrm{eff} \leq \|\hatH_\CM\|_\mathrm{op}/\hbar_\mathrm{eff}$.  Summing gives~\eqref{eq:centbound}.  Equality is attained since $S(\id/d) = \log d$ and $\Tr((\id/d)\hatH_\CM) = \Tr(\hatH_\CM)/d$.
\end{proof}

\begin{theorem}[Constraint Capacity Theorem]
\label{thm:capacity}
The maximum mutual information between the reality state and the cognitive
state, called the \emph{constraint channel capacity}, is bounded by:
\begin{equation}
  C_\Phi \leq \log\dim\HC_\CM
    - \max_{\sigma \in \mathfrak{D}(\CR_\mathrm{op})} S\bigl(\Phi_X(\sigma)\bigr)
    + S(\Phi_X(\sigma_\mathrm{max})),
  \label{eq:capacity}
\end{equation}
where $\sigma_\mathrm{max}$ achieves the maximum output entropy.
\end{theorem}

\begin{proof}
Standard quantum channel capacity bound (Holevo bound \cite{Holevo1973}):
$C_\Phi \leq \chi(\{p_i,\sigma_i\}) = S(\sum_i p_i\Phi(\sigma_i)) - \sum_i p_i S(\Phi(\sigma_i))$;
optimising over input ensembles and using $\dim\HC_\CM$ as the output space dimension gives~\eqref{eq:capacity}.
\end{proof}

\begin{theorem}[Cognitive Superposition Principle]
\label{thm:superposition}
If $\psi_1, \psi_2 \in \HC_\CM$ are valid cognitive states, then any
normalised linear combination $\alpha\psi_1 + \beta\psi_2$
($|\alpha|^2+|\beta|^2=1$) is also a valid cognitive state in $\HC_\CM$.
\end{theorem}

\begin{proof}
$\HC_\CM$ is a Hilbert subspace by Definition~\ref{def:subspace} and hence closed under linear combination.  The normalised combination lies in $\HC_\CM$ since $\hatP_\CM(\alpha\psi_1+\beta\psi_2) = \alpha\hatP_\CM\psi_1 + \beta\hatP_\CM\psi_2 = \alpha\psi_1+\beta\psi_2$.
\end{proof}

\begin{theorem}[No-Cloning in CIIR]
\label{thm:nocloning}
There is no CPTP map $\Lambda : \CB(\HC) \to \CB(\HC\otimes\HC)$ such that
$\Lambda(\ket\psi\bra\psi) = \ket\psi\bra\psi \otimes \ket\psi\bra\psi$
for all $\ket\psi \in \HC_\CM$.
\end{theorem}

\begin{proof}
The proof is identical to the standard quantum no-cloning theorem
\cite{Wootters1982}: suppose such $\Lambda$ exists.  For two
non-orthogonal states $\psi_1,\psi_2$ with $\langle\psi_1|\psi_2\rangle = c \neq 0$,
unitarity of $\Lambda$'s Stinespring dilation gives
$c = \langle\psi_1|\psi_2\rangle^2$, which is impossible for $c \in (0,1)$.
\end{proof}

\begin{theorem}[Interface Continuity]
\label{thm:continuity}
The interface map $\Phi_X : \mathfrak{D}(\CR_\mathrm{op}) \to \mathfrak{D}(\HC)$
is uniformly continuous with respect to the trace norm:
\begin{equation}
  \|\Phi_X(\sigma_1) - \Phi_X(\sigma_2)\|_1
    \leq \|\sigma_1 - \sigma_2\|_1
  \label{eq:continuity}
\end{equation}
for all $\sigma_1, \sigma_2 \in \mathfrak{D}(\CR_\mathrm{op})$.
\end{theorem}

\begin{proof}
Every CPTP map is a contraction under the trace norm: this is a standard
result following from data processing inequality for quantum channels
\cite{Lindblad1975}.  Hence $\|\Phi_X(\sigma_1)-\Phi_X(\sigma_2)\|_1
\leq \|\sigma_1-\sigma_2\|_1$.
\end{proof}

% ============================================================
\section{Model Construction}
\label{sec:models}
% ============================================================

\subsection{Finite-Dimensional Model}
\label{subsec:finitedim}

\begin{example}[Qubit Constraint Model]
\label{ex:qubit}
Let $\HC = \mathbb{C}^2$ (single qubit), $\CM = S^1$ (unit circle),
and $\hatH_\CM = \omega\hat\sigma_z$ where $\hat\sigma_z$ is the Pauli-$Z$
operator and $\omega > 0$ is a constraint frequency.  The interface map is
\[
  \Phi_X(\theta) = \cos^2(\theta/2)\ket{0}\bra{0} + \sin^2(\theta/2)\ket{1}\bra{1},
  \quad \theta \in [0,2\pi),
\]
which is CPTP and surjective onto diagonal states in $\mathfrak{D}(\mathbb{C}^2)$.
The CIIR master equation~\eqref{eq:master} reduces to:
\[
  \dot\rho = -i\omega[\hat\sigma_z,\rho] + \gamma({\hat\sigma_-}\rho{\hat\sigma_+}
    - \tfrac12\hat\sigma_+\hat\sigma_-\rho - \tfrac12\rho\hat\sigma_+\hat\sigma_-),
\]
the standard two-level Lindblad equation.  This model satisfies all six axioms
and hence $\mathcal{M}_\mathrm{qubit} \models \mathcal{A}$.
\end{example}

\begin{theorem}[Finite-Dimensional Model Existence]
\label{thm:finitedim}
For any $n \in \mathbb{N}$ and constraint manifold $\CM$ with
$\kappa_\mathrm{max}\cdot\mathrm{vol}(\CM) \leq \log n$, there exists a valid
CIIR model with $\HC = \mathbb{C}^n$ satisfying all six axioms.
\end{theorem}

\begin{proof}
Take $\HC = \mathbb{C}^n$, $\hatH_\CM = \mathrm{diag}(\epsilon_1,\ldots,\epsilon_n)$
with $\epsilon_k = k\hbar_\mathrm{eff}\kappa_\mathrm{max}/n$, and
$\Phi_X = $ the pinching map onto diagonal states.  All axioms are satisfied:
Axiom~\ref{ax:hilbert} by $\dim\HC = n \leq \exp(\kappa_\mathrm{max}\cdot\mathrm{vol}(\CM))$;
the others by direct verification.
\end{proof}

\subsection{Infinite-Dimensional Model}
\label{subsec:infinitedim}

\begin{example}[Harmonic Oscillator Constraint Model]
\label{ex:oscillator}
Let $\HC = L^2(\mathbb{R})$, $\CM = \mathbb{R}_+$,
$\hatH_\CM = \hbar_\mathrm{eff}\hat\omega(\hat{a}^\dagger\hat{a} + \tfrac12)$
(quantum harmonic oscillator), and interface map
\[
  \Phi_X(f) = \sum_{n=0}^\infty |\langle n|f\rangle|^2 \ket{n}\bra{n}
\]
for $f \in L^2(\mathbb{R})$ with $\|f\|=1$.  The cognitive state
$\rho(t) = e^{t\mathcal{L}}\rho(0)$ with Lindblad operators
$L_1 = \sqrt{\gamma}\hat{a}$, $L_2 = \sqrt{\gamma n_\mathrm{th}}\hat{a}^\dagger$
gives thermal relaxation to $\rho_* = (1-e^{-\beta\hbar_\mathrm{eff}\hat\omega})\sum_n e^{-n\beta\hbar_\mathrm{eff}\hat\omega}\ket{n}\bra{n}$.
This satisfies all six axioms and serves as $\mathcal{M}_\mathrm{osc} \models \mathcal{A}$.
\end{example}

\begin{theorem}[Infinite-Dimensional Model Existence]
\label{thm:infinitedim}
For $\CM = [0,\infty)$ with constant curvature $\kappa_0 > 0$ and
$\mathrm{vol}(\CM) = \infty$, the CIIR system with $\HC = L^2(\mathbb{R})$
satisfies all six axioms.
\end{theorem}

\begin{proof}
Axiom~\ref{ax:hilbert} permits $\dim(\HC) = \infty$ when
$\kappa_\mathrm{max}\cdot\mathrm{vol}(\CM) = \infty$.  The harmonic oscillator
Hamiltonian is self-adjoint on $L^2(\mathbb{R})$ with dense domain
$\mathscr{S}(\mathbb{R})$ (Schwartz space).  The remaining axioms follow by the
same verification as in Example~\ref{ex:oscillator}.
\end{proof}

% ============================================================
\section{Empirical Mapping}
\label{sec:empirical}
% ============================================================

\begin{definition}[Empirical Interpretation Map]
\label{def:empmap}
The \emph{empirical interpretation map} is a measurable function
\begin{equation}
  \CI : \mathfrak{D}(\HC) \to \mathbb{R}^n,
  \quad
  \CI(\rho) = \bigl(\Tr(\rho\hat{O}_1),\, \Tr(\rho\hat{O}_2),\,\ldots,\,
                    \Tr(\rho\hat{O}_n)\bigr),
  \label{eq:empmap}
\end{equation}
where $\{\hat{O}_1,\ldots,\hat{O}_n\} \subset \CA(\HC)$ is a tomographically
complete set of observables.
\end{definition}

\begin{proposition}[Tomographic Completeness]
\label{prop:tomo}
A set $\{\hat{O}_k\}$ is tomographically complete if and only if
$\{\hat{O}_k\}$ spans $\CB(\HC)$ as a real vector space; in this case $\CI$
is injective on $\mathfrak{D}(\HC)$.
\end{proposition}

\begin{proof}
Standard result: density operators are uniquely determined by their expectation
values on a spanning set of observables \cite{Prugovecki1977}.
\end{proof}

The empirical map $\CI$ connects:
\begin{itemize}
\item \textit{Measurable cognitive outputs}: $\CI_k(\rho) = \Tr(\rho\hat{O}_k)$
      represents the mean response on behavioural task $k$.
\item \textit{Behavioral probabilities}: For dichotomic observable
      $\hat{O} = \hatP_+ - \hatP_-$, $\CI(\rho) = \Tr(\hatP_+\rho)$ is the
      probability of the positive outcome.
\item \textit{Information-theoretic quantities}: Taking
      $\hat{O} = -\log\rho_*$ gives $\CI(\rho) = S(\rho\|\rho_*)$, the
      relative entropy with respect to the equilibrium.
\end{itemize}

% ============================================================
\section{Predictive Framework and Falsifiable Predictions}
\label{sec:predictions}
% ============================================================

\begin{theorem}[Constraint-Bandwidth Trade-off]
\label{thm:bandwidth}
For a cognitive system with constraint manifold $\CM$ of volume $V$, the number
of distinguishable interface states (channel capacity in nats) satisfies:
\begin{equation}
  C_\Phi \leq \kappa_\mathrm{max} \cdot V + \log\dim\HC,
  \label{eq:bandwidth}
\end{equation}
with saturation iff the Lindblad operators vanish ($\CD_\CM = 0$).
\end{theorem}

\begin{proof}
From Theorem~\ref{thm:capacity} and $\dim\HC_\CM \leq e^{\kappa_\mathrm{max}V}$,
we obtain $C_\Phi \leq \log\dim\HC_\CM + \log\dim\HC \leq \kappa_\mathrm{max}V + \log\dim\HC$.
Saturation requires zero decoherence.
\end{proof}

\subsection{Quantitative Falsifiable Predictions}

\textbf{Prediction 1: Curvature-Induced Decision Latency.}
In a cognitive task where the constraint manifold curvature is parametrically
varied (e.g.\ via task complexity $\kappa$ measured as Kolmogorov complexity of
the stimulus), the mean decision latency $\tau$ satisfies:
\begin{equation}
  \tau = \tau_0 \exp\!\bigl(\alpha\,\kappa \cdot V_\CM\bigr) + \tau_1,
  \label{eq:pred1}
  \tag{P1}
\end{equation}
with $\alpha > 0$ a system-specific constant, $\tau_0, \tau_1 > 0$.
\emph{Falsification:} If $\tau$ grows sub-exponentially in $\kappa V_\CM$, or
if $\tau$ is independent of $V_\CM$ at fixed $\kappa$, Prediction~1 is
falsified.

\textbf{Prediction 2: Decoherence-Scaling of Cognitive Interference.}
When two cognitive systems $X_1,X_2$ with identical constraint manifolds
interact through tensor-product coupling, the off-diagonal interference term
in Theorem~\ref{thm:interference} decays as:
\begin{equation}
  |\mathrm{Re}\langle\psi_1|\hatO|\psi_2\rangle(t)| =
    |\mathrm{Re}\langle\psi_1|\hatO|\psi_2\rangle(0)|\,
    e^{-(\Gamma_1+\Gamma_2)t/2},
  \label{eq:pred2}
  \tag{P2}
\end{equation}
where $\Gamma_k = \sum_j \gamma_{j}^{(k)} \|L_j^{(k)}\|^2$ are the total
decoherence rates.  \emph{Falsification:} Super-exponential decay, oscillatory
recovery, or independence from $\Gamma_k$ would falsify~P2.

\textbf{Prediction 3: Constraint Capacity Saturation.}
The mutual information between environmental stimuli and cognitive responses in
a well-trained agent saturates at:
\begin{equation}
  I(R;C) \to C_\Phi = \kappa_\mathrm{max}\cdot V_\CM + O(1)
    \quad \text{as } T \to \infty,
  \label{eq:pred3}
  \tag{P3}
\end{equation}
where $T$ is training time and $V_\CM$ is the learnable constraint volume.
\emph{Falsification:} If $I(R;C)$ fails to saturate, or saturates at a
value exceeding $\kappa_\mathrm{max}V_\CM + O(1)$, prediction~P3 is falsified.

% ============================================================
\section{Comparative Embedding}
\label{sec:embedding}
% ============================================================

\begin{theorem}[Classical Probability as Limiting Case]
\label{thm:classicallimit}
When $\CM$ is flat ($\kappa = 0$, $K=0$), the observable algebra
$\CA(\HC)$ is Abelian (all observables commute), and CIIR reduces to
classical Kolmogorov probability theory on the spectrum
$\mathrm{spec}(\CA)$.
\end{theorem}

\begin{proof}
By Corollary~\ref{cor:flat}, $\kappa = 0 \Rightarrow [\hatC_i,\hatC_j]=0$ for all $i,j$.  The full observable algebra is then Abelian, and by the spectral theorem for commuting self-adjoint operators, $\CA(\HC)$ is $*$-isomorphic to $L^\infty(\Omega, \mu)$ for some measure space $(\Omega,\mu)$.  Density operators on $L^\infty(\Omega,\mu)$ correspond to probability measures on $\Omega$, recovering Kolmogorov's framework.
\end{proof}

\begin{theorem}[Bayesian Inference as Measurement Update]
\label{thm:bayesian}
In the flat, Abelian regime, the post-measurement update rule~\eqref{eq:collapse}
reduces to Bayes' theorem:
\begin{equation}
  P(H|E) = \frac{P(E|H)P(H)}{P(E)}.
  \label{eq:bayes}
\end{equation}
\end{theorem}

\begin{proof}
In the Abelian case $\rho \leftrightarrow P$ (probability distribution), $\hat{P}_E \leftrightarrow \mathbb{1}_E$ (indicator function).  The L{\"u}ders rule~\eqref{eq:collapse} gives $P'(H) = P(H \cap E)/P(E) = P(E|H)P(H)/P(E)$, which is Bayes' theorem.
\end{proof}

\begin{theorem}[Quantum Cognition as Special Case]
\label{thm:qcog}
Quantum cognition models (e.g.\ Busemeyer--Bruza \cite{Busemeyer2012}) are
special cases of CIIR in which $\CM$ is a flat manifold and
$\HC = \mathbb{C}^n$ for small $n$, with $\Phi_X$ the identity map and
$\hatH_\CM$ a context-encoding Hamiltonian.
\end{theorem}

\begin{proof}
Quantum cognition posits a Hilbert space $\mathbb{C}^n$ with projective
measurements \cite{Busemeyer2012}.  Setting $\CM = \{*\}$ (a single-point
flat manifold), $\kappa=0$, $\Phi_X = \id$, and $\hatH_\CM = 0$ recovers
the basic quantum cognition setup.  Context effects are then encoded by
non-trivial $\hatH_\CM$ in the CIIR framework, providing a principled
derivation of context operators absent from standard quantum cognition.
\end{proof}

% ============================================================
\section{Discussion}
\label{sec:discussion}
% ============================================================

\subsection{Relationship to Information-Theoretic Interpretations}

CIIR is compatible with relational quantum mechanics \cite{Rovelli1996} and
QBism \cite{Fuchs2014} in that observables are indexed by agents, but it
differs in grounding interface projections in differential geometry rather
than subjective Bayesian beliefs.  The constraint manifold $\CM$ is
a structural feature of the system, not a belief state.

\subsection{Connection to Predictive Coding}

The CIIR master equation~\eqref{eq:master} is structurally analogous to
free-energy minimisation in predictive coding \cite{Friston2010}: the
Hamiltonian term drives the state toward a prediction and the dissipator
models prediction-error updating.  A formal equivalence holds when
$\hatH_\CM$ is the negative log-evidence (surprise) operator.

\subsection{Computational Complexity}

Simulating an $n$-qubit CIIR system requires storing $\rho \in \mathfrak{D}(\mathbb{C}^{2^n})$, which is exponential in $n$.  This is consistent with the Axiom~\ref{ax:hilbert} bound $\dim\HC = O(\exp(\kappa V))$ and confirms that realistic cognitive systems must operate with bounded $\kappa V$.

% ============================================================
\section{Limitations}
\label{sec:limitations}
% ============================================================

\begin{enumerate}
\item \textbf{Markov approximation:} The GKSL master equation assumes
      memoryless environment coupling; non-Markovian extensions require
      Nakajima--Zwanzig projector techniques.
\item \textbf{Constraint manifold specification:} CIIR does not uniquely
      determine $\CM$ from observable data; additional empirical input is
      required.
\item \textbf{Ontological commitments:} Axiom~\ref{ax:ontological} commits
      to a mind-independent $\CR$ whose metric structure is not further
      constrained by the present theory.
\item \textbf{Classical limit uniqueness:} Theorem~\ref{thm:classicallimit}
      shows flatness is \emph{sufficient} for classicality; uniqueness of
      the classical limit requires additional decoherence arguments.
\end{enumerate}

% ============================================================
\section{Conclusion}
\label{sec:conclusion}
% ============================================================

We have constructed a complete, closed, axiomatic theory $\TCIIR =
(\CS,\CA,\CR,\CD,\CM,\CI)$ for Constraint-Induced Interface Realism.
Starting from six independent axioms, we have:
\begin{enumerate}
\item Defined all primitive objects (reality space, constraint manifold,
      cognitive state space, interface map, observable algebra).
\item Proved the CIIR Representation Theorem establishing functorial
      Hilbert-space representations.
\item Derived the CIIR master equation and its stability properties.
\item Established ten non-trivial theorems with formal proofs.
\item Constructed finite- and infinite-dimensional model classes and verified
      $\mathcal{M} \models \mathcal{A}$.
\item Derived three quantitatively falsifiable predictions.
\item Embedded classical probability, Bayesian inference, and quantum
      cognition as special cases.
\end{enumerate}

The framework is sufficiently rigorous that a mathematician or theoretical
physicist can directly extend, critique, or attempt to falsify it from the
formulation given.

% ============================================================
\appendix
% ============================================================

\section{Derivation of Constraint Commutation Relations}
\label{app:commutation}

Here we give an expanded derivation of Proposition~\ref{prop:commutation}.
Let $\CM$ be a Riemannian manifold with connection $\nabla$.  The curvature
2-form $F$ is defined by $F(u,v)w = \nabla_u\nabla_v w - \nabla_v\nabla_u w - \nabla_{[u,v]}w$ for vector fields $u,v,w$.  The constraint operator $\hatC$ is a smooth field of bounded operators on $\HC$, and its parallel transport is defined via the pullback connection $\Phi_X^*\nabla$.  The computation in Section~\ref{sec:algebra} then follows from the holonomy representation of curvature.

\section{Proof of Entropy Non-Decrease}
\label{app:entropy}

We expand the proof of Theorem~\ref{thm:entropy} for the general non-primitive case.
Define $f(t) = S(\rho(t)\|\rho_*)$.  By the data processing inequality, $f(t) \leq f(0)$ for all $t \geq 0$ under any CPTP semigroup.  Hence $f$ is non-increasing.  Separating $f = -S(\rho) - \Tr(\rho\log\rho_*)$ and noting $\Tr(\rho\log\rho_*)$ may increase, the net effect on $S(\rho)$ depends on the equilibrium structure.  For a primitive dissipator with unique $\rho_*$, $-\Tr(\rho\log\rho_*)$ converges to a fixed constant and $S(\rho)$ must increase to compensate the decrease in $f$.

\section{Qubit Lindblad Explicit Form}
\label{app:qubit}

For the qubit model (Example~\ref{ex:qubit}):
\[
  \rho = \begin{pmatrix}\rho_{00} & \rho_{01}\\ \bar\rho_{01} & \rho_{11}\end{pmatrix},
\]
the master equation gives:
\begin{align}
  \dot\rho_{00} &= -\gamma\rho_{00} + \gamma\rho_{11}, \\
  \dot\rho_{11} &= \gamma\rho_{00} - \gamma\rho_{11}, \\
  \dot\rho_{01} &= -(i\omega + \gamma/2)\rho_{01}.
\end{align}
The equilibrium is $\rho_* = \frac12\id$ and $\Gamma_{01} = \gamma/2$, consistent with Theorem~\ref{thm:decoherence}.

\bibliography{references}

\end{document}
